% =================================================
% HOVEDDOKUMENT (KLASSISK ARTIKKEL)
% =================================================
%
% Klassisk LaTeX-artikkel. Layout og formatering ligger
% i klassiskstil.sty. Font, matematikkpakker og makroer
% velges her.
%
% =================================================


\documentclass[12pt]{article}


% =================================================
% LAST INN STIL
%
% Verktøy tilgjengelig via klassiskstil.sty:
%   - figurer:    graphicx, float ([H]-plassering), wrapfig, subfig
%   - tabeller:   booktabs (\toprule, \midrule, \bottomrule), caption
%   - lister:     enumitem (egendefinerte etiketter, avstander)
%   - bokser:     tcolorbox (fargede ramme- og notatbokser)
%   - farger:     xcolor med svgnames
%   - siteringer: natbib (numeriske, sortert og komprimert)
%
% Se README.md for syntakseksempler.
% =================================================
\usepackage{klassiskstil}


% =================================================
% SPRÅK
% =================================================
\usepackage[norsk]{babel}


% =================================================
% FONT
%
% Velg én tekstfont. Kommentér ut de andre.
% =================================================

\usepackage{mlmodern}              % Modern Latin Modern
% \usepackage{lmodern}             % Latin Modern (klassisk)
% \usepackage{stix2}               % STIX2 (matte og tekst)

% Sansserif-font (valgfritt):
\usepackage{helvet}                % Helvetica som sans-font

% For sansserif-overskrifter, fjern kommentartegnene:
% \usepackage{titlesec}
% \titleformat*{\section}{\sffamily\Large\bfseries}
% \titleformat*{\subsection}{\sffamily\large\bfseries}


% =================================================
% MATEMATIKK OG VITENSKAP
% =================================================
\usepackage{amsmath, amsfonts, amsthm}
\usepackage{sansmath}              % matematikk i sans
\usepackage{siunitx}               % SI-enheter
\usepackage{esvect}                % fine vektorer


% =================================================
% EGENDEFINERTE PAKKER
%
% Legg til prosjektspesifikke pakker her.
% =================================================

% \usepackage{tikz}                % tegninger / diagrammer
% \usepackage{physics}             % \dd, \ket, \bra, ...


% =================================================
% EGENDEFINERTE KOMMANDOER
%
% Personlige matematiske snarveier.
% =================================================

% --- Greske bokstaver / konstanter ---
\newcommand{\p}{\varphi}
\newcommand{\ep}{\varepsilon}
\newcommand{\e}{\mathrm{e}}

% --- Enhetsvektorer ---
\newcommand{\er}{\vec{e}_r}
\newcommand{\ephi}{\vec{e}_\varphi}

% --- Derivasjon ---
\newcommand{\pdd}[2]{\frac{\partial #1}{\partial #2}}
\newcommand{\pddd}[1]{\frac{\partial^2}{\partial #1^2}}
\newcommand{\dd}[2]{\frac{\mathrm{d} #1}{\mathrm{d} #2}}
\newcommand{\diff}[1]{\mathrm{d}{#1}~}
\newcommand{\diffn}[2]{\mathrm{d}^{#1} {#2}~}
\newcommand{\diffnn}[2]{\mathrm{d}^{#2} {#1}~}

% --- Partielle deriverte ---
\newcommand{\ddx}{\frac{\partial }{\partial x}}
\newcommand{\ddy}{\frac{\partial }{\partial y}}
\newcommand{\ddr}{\frac{\partial }{\partial r}}
\newcommand{\dddx}{\frac{\partial^2}{\partial x^2}}
\newcommand{\dddy}{\frac{\partial^2}{\partial y^2}}
\newcommand{\ddphi}{\frac{\partial }{\partial \varphi}}
\newcommand{\dddr}{\frac{\partial^2 }{\partial r^2}}
\newcommand{\dddphi}{\frac{\partial^2 }{\partial \varphi^2}}

% --- Integrasjon ---
\newcommand{\pint}[2]{\int\!{#1}\,d{#2}}
\newcommand{\pintlim}[4]{\int_{#1}^{#2}\!{#3}\,d{#4}}

% --- Algebra ---
\newcommand{\nsqrt}[2]{\sqrt[\uproot{2}\leftroot{5} {#1}] {#2}}


% =================================================
% METADATA
% =================================================
\title{\huge Tittel på artikkelen}
\author{\small Fornavn Etternavn}
\date{Våren 2026}


% =================================================
% DOKUMENT
% =================================================
\begin{document}

{\fontfamily{lmr}\selectfont\maketitle}     % tittel i Latin Modern
\vspace{1cm}

\begin{quote}
  \emph{``Et klassisk sitat plasseres her, og setter tonen for artikkelen
  som følger.''}
  \begin{flushright}
    --- Forfatter
  \end{flushright}
\end{quote}

\begin{abstract}
  \noindent Sammendraget gir leseren et raskt overblikk over hva
  artikkelen handler om, hva som er det sentrale argumentet, og hvorfor
  det er relevant. To til fire setninger er vanligvis tilstrekkelig.
\end{abstract}

\newpage


\section{Innledning}

Innledningen introduserer problemstillingen og setter den i kontekst.
Avsnitt markeres med innrykk og ikke med linjeskift, slik at teksten
danner en sammenhengende strøm. Dette er den klassiske typografiske
konvensjonen for trykte artikler og bøker.

Et nytt avsnitt fortsetter med samme rytme. Når det refereres til
tidligere arbeid, gjøres dette med en sitering~\cite{eksempel2024}, og
flere kilder kan kombineres i samme parentes~\cite{eksempel2024,eksempelbok2020}.


\section{Sentrale resultater}

Matematikk settes enten inline, som $E = mc^2$, eller som en nummerert
ligning:
\begin{equation}
  \hat{H}\psi = E\psi.
  \label{eq:eksempel}
\end{equation}
Det er naturlig å henvise til ligninger med \eqref{eq:eksempel} der det
passer i teksten.

Et nytt avsnitt etter en ligning starter med innrykk akkurat som
ellers. Avstanden mellom ligning og brødtekst er satt til 0,5 cm i
\texttt{klassiskstil.sty}, slik at matematikken puster uten å løsrive
seg fra resonnementet.

\begin{figure}[H]
  \centering
  \includegraphics[width=0.7\linewidth]{Energy_levels}
  \caption{Et eksempel på en figur. \texttt{[H]}-plasseringen låser
           figuren akkurat der den står i kilden, mens \texttt{[t]}
           lar den flyte til toppen av siden.}
  \label{fig:eksempel}
\end{figure}

Henvis til figurer i teksten med figur~\ref{fig:eksempel}. For å
oppsummere nøkkelverdier kan tabeller være praktiske:

\begin{table}[H]
  \centering
  \caption{Eksempel på en tabell satt med \texttt{booktabs}.}
  \label{tab:eksempel}
  \begin{tabular}{lcr}
    \toprule
    Størrelse           & Symbol  & Verdi                         \\
    \midrule
    Lysets hastighet    & $c$     & $2{,}998 \times 10^8$ m/s     \\
    Plancks konstant    & $\hbar$ & $1{,}055 \times 10^{-34}$ Js  \\
    Elementærladning    & $e$     & $1{,}602 \times 10^{-19}$ C   \\
    \bottomrule
  \end{tabular}
\end{table}

Vertikale streker er bevisst utelatt, og horisontale brukes sparsomt
--- bare topp, midt og bunn. Det gir tabellen et rolig, profesjonelt
preg.


\section{Diskusjon}

Diskusjonen er stedet hvor man tolker resultatene og setter dem i
sammenheng med eksisterende litteratur. Her kan man være mer
spekulativ enn i resultatdelen, så lenge det er tydelig hva som er
egne resonnementer og hva som er etablerte funn.

Punktlister er nyttige når flere likestilte poeng skal listes opp:
\begin{itemize}
  \item Det første poenget formuleres her.
  \item Et nytt poeng følger på neste linje.
  \item Et tredje poeng avslutter listen.
\end{itemize}

For nummererte sekvenser eller stegvise resonnement bruker man
nummererte lister:
\begin{enumerate}
  \item Første steg i resonnementet.
  \item Andre steg, som bygger på det første.
  \item Tredje steg leder fram til konklusjonen.
\end{enumerate}

\begin{tcolorbox}[colback=Beige, colframe=DarkSlateGray, title=Merknad]
  Fargede bokser er nyttige for å fremheve definisjoner, observasjoner
  eller advarsler. Pakken \texttt{tcolorbox} gir mange muligheter for
  utforming --- se dokumentasjonen for fulle valg.
\end{tcolorbox}

\subsection{En underseksjon}

Underseksjoner får sin egen overskrift, men beholder en rolig
typografisk vekt slik at hovedseksjonene fortsatt dominerer
hierarkiet.


\section{Avslutning}

Avslutningen oppsummerer kort, peker mot åpne spørsmål, og foreslår
videre arbeid hvis det er relevant. Et siste avsnitt avrunder
artikkelen uten å introdusere nye argumenter.


% =================================================
% REFERANSER
% Kompileringsrekkefølge:
%   pdflatex main
%   bibtex   main
%   pdflatex main
%   pdflatex main
% =================================================
\bibliographystyle{unsrtnat}
\bibliography{referanser}


\end{document}
